\subsection{Conclusión individual: Belinda Toaquiza}

El desarrollo de la base compartida y la comunicación TCP me permitió relacionar
los conceptos de sistemas distribuidos con una implementación de tres nodos.
La configuración común y los identificadores de mensajes facilitan la integración
del trabajo y el seguimiento de cada solicitud y su acuse. También comprendí la
importancia de delimitar los mensajes JSON, proteger el reloj ante accesos
concurrentes y conservar registros que permitan revisar lo ocurrido.

Las pruebas y validaciones locales completaron los intercambios sin fallos y
permitieron comprobar que el reloj de cada recepción fuera mayor que el de su
envío correspondiente. La comparativa final del equipo registró menor latencia
media y mayor rendimiento en TCP para las cuatro combinaciones evaluadas.
Estos resultados corresponden a las condiciones medidas. Ejecutar los nodos
en una sola máquina también limita
las conclusiones sobre una red entre equipos distintos. Como aprendizaje,
destaco que documentar las condiciones, registrar los errores y conservar los
datos originales permite revisar y reproducir el trabajo.
